\section{Robustnost při misspecifikaci prioru}
\label{sec:misspecifikace}

V této poslední sekci budeme zkoumat vlastnost, na které selhává většina Bayesovských metod, a tou je tzv. \emph{misspecifikace prioru}.
Situaci, kdy jsme prior apriorně odhadli špatně, tedy když skutečná data pocházejí z jiného rozdělení, než jaké prior předpokládá, nazýváme misspecifikací a tento pojem budeme v celé sekci používat.
Metodě s takto špatně zvoleným priorem pak budeme říkat misspecifikovaná.
Ve všech dosavadních experimentech pocházela data z GP s RBF kernelem, tedy přesně z té třídy modelů, kterou porovnávané metody předpokládaly.
Vždy proto existoval správně nastavený GP (\emph{oracle}), ke kterému jsme mohli PFN vztahovat.
V praxi ale taková jistota neexistuje, skutečná data nepřicházejí se štítkem, ze kterého kernelu byla vygenerována.
Nyní proto odstraníme i tuto poslední jistotu.
Data budeme generovat z tzv. \emph{Matérnova kernelu} (definice viz dále), zatímco všechny porovnávané metody budou předpokládat RBF.
Skutečný generativní proces tak bude ležet mimo support prioru všech metod, včetně implicitního prioru PFN.
Budeme tedy zkoumat, co se stane, když po PFN chceme predikovat mimo support jeho prioru, o kolik se predikce zhorší a která metoda za těchto podmínek selže nejméně.
Připomeňme, že přesně touto situací se zabývala věta~\ref{theorem:PPDcanLearn}, podle které se PPD umí učit i tehdy, když skutečné rozdělení leží mimo support prioru.
Tato sekce je tedy zároveň empirickým testem této teoretické vlastnosti.
\par

Nejprve popíšeme uspořádání experimentu.
Data generujeme z GP s \emph{Matérnovým kernelem} s korelační délkou $\ell_{\text{true}}$ a šumem $\sigma^2 = 0{,}1$.
Matérnův kernel je v praxi nejběžnější alternativou RBF a jeho tvar řídí parametr hladkosti $\nu$.
My volíme $\nu = 2{,}5$, pro který má kernel tvar
$$k_{\text{Matérn}}(x, x') = \left(1 + \frac{\sqrt{5}\,r}{\ell} + \frac{5 r^2}{3 \ell^2}\right) \exp\!\left(-\frac{\sqrt{5}\,r}{\ell}\right), \qquad r = |x - x'|,$$
kde $\ell$ je opět korelační délka.
Realizace GP s tímto kernelem jsou jen dvakrát diferencovatelné, tedy znatelně hrubší než nekonečně hladké realizace RBF kernelu.
Právě proto se Matérnův kernel považuje za realističtější model skutečných dat, nekonečná hladkost RBF je totiž pro reálné procesy příliš silný předpoklad [Rasmussen].
Náš test misspecifikace tak napodobuje typickou praktickou situaci, kdy jsou data hrubší, než model předpokládá.
Pro tento účel je Matérnův kernel ideální i z druhého důvodu.
Pro $\nu \to \infty$ přechází v RBF, oba kernely jsou si tedy dost podobné na to, aby jejich záměna byla uvěřitelná, a zároveň dost odlišné na to, aby ji žádná volba hyperparametrů RBF nedokázala úplně napravit.
Misspecifikace prioru zde tedy spočívá v rozdílu mezi skutečným a předpokládaným kernelem.
Žádná z testovaných metod nemá k dispozici správnou třídu funkcí.
Porovnáváme následující metody:
\begin{itemize}
    \item \emph{PFN}, náš šestivrstvý model s náhodnými hyperparametry ze sekce~\ref{sec:pfn-random-hp}, jehož prior je implicitně daný trénovacím rozdělením (RBF, $\ell \sim \mathcal{U}(0{,}05,\, 1{,}0)$) a za běhu ho nelze nijak opravit,
    \item \emph{GP\_fixed}, GP s RBF kernelem a pevnou korelační délkou $\ell = 0{,}7$, tedy špatný kernel i špatná korelační délka a žádná adaptace,
    \item \emph{GP\_ML}, GP s RBF kernelem, který korelační délku odhaduje z dat maximalizací marginální věrohodnosti (Type II – MLE ze sekce~\ref{subsec:type2-mle}),
    \item \emph{GP\_Bayes}, Bayesovský GP s RBF kernelem a priorem $p(\ell) \sim \text{Gamma}(4,\, 4)$, vzorkovaný metodou NUTS ze sekce~\ref{subsec:mcmc},
    \item \emph{GP\_oracle}, referenční GP se správným Matérnovým kernelem i správnými hyperparametry.
\end{itemize}
Poznamenejme ještě, jak je nastavený prior metody GP\_Bayes.
Rozdělení $\text{Gamma}(4,\, 4)$ má střední hodnotu u $\ell = 1{,}0$ a většinu své pravděpodobnostní hmoty soustřeďuje kolem ní.
Pravá hodnota $\ell_{\text{true}} = 0{,}3$ leží až v jeho levém chvostu, jak ukazuje obrázek~\ref{fig:exp7_prior}.
Prior tedy pravé hodnotě přiřazuje jen malou pravděpodobnost.
GP\_Bayes je tím pádem misspecifikovaný dvakrát.
Za prvé předpokládá špatný kernel, stejně jako všechny ostatní metody.
Za druhé jeho prior tvrdí, že správná korelační délka je nepravděpodobná, a metoda se od tohoto přesvědčení může odchýlit až na základě pozorovaných dat.
Experiment se tím blíží reálné situaci, kdy prior volíme bez znalosti pravdy a zvolíme ho nešťastně.
\par

\begin{figure}[htbp]
    \centering
    \includegraphics[width=0.82\linewidth]{figures/GP2_exp7_misspecification/fig_exp7_prior_visualization.png}
    \caption{Prior $p(\ell) \sim \text{Gamma}(4,\, 4)$ metody GP\_Bayes (fialová) se střední hodnotou u $\ell = 1{,}0$. Zelená přerušovaná čára značí pravou hodnotu $\ell_{\text{true}} = 0{,}3$, která leží v chvostu prioru. Modrý pás vyznačuje trénovací rozsah PFN $\ell \in [0{,}05,\, 1{,}0]$, pravá hodnota leží uvnitř něj.}
    \label{fig:exp7_prior}
\end{figure}

Kvalitu predikce měříme dvěma metrikami.
První je MSE predikce vůči skutečným hodnotám.
Druhou je kalibrační metrika, které budeme říkat \emph{$\Delta$NLL}:
$$\Delta\text{NLL} = \text{NLL}_{\text{metoda}} - \text{NLL}_{\text{oracle}},$$
tedy o kolik je záporná log-věrohodnost predikcí dané metody horší než u \emph{oracle}.
\emph{Oracle} má z definice $\Delta\text{NLL} = 0$ a čím blíže nule metoda je, tím lépe je kalibrovaná.
Vysvětleme, v čem se obě metriky liší a proč potřebujeme obě.
MSE hodnotí pouze bodovou predikci, tedy střed prediktivního rozdělení.
Dvě metody se stejným středem, ale různě širokou nejistotou, mají proto stejné MSE.
NLL naproti tomu hodnotí celé prediktivní rozdělení.
Trestá nejen špatný střed, ale i špatně odhadnutou nejistotu, a to zejména příliš sebejisté intervaly, ze kterých skutečná hodnota vypadne.
Pořadí metod se tak v obou metrikách může lišit, metoda s dobrým středem a přehnaně úzkou nejistotou dopadne v MSE dobře a v $\Delta$NLL špatně.
Právě $\Delta$NLL přitom testuje to podstatné, smyslem Bayesovské inference je totiž vracet celou prediktivní nejistotu, nikoli jen bodový odhad.
MSE pak slouží jako kontrola, že případná lepší kalibrace není zaplacená horší bodovou predikcí.
Všechny výsledky průměrujeme přes $150$ instancí, v každé predikujeme $n_{\text{test}} = 10$ testovacích bodů.
\par

Než přejdeme k samotným testům, ujasněme si, co přesně zde misspecifikujeme.
Priorem každé metody rozumíme její představu o tom, z jaké třídy funkcí data pocházejí.
U GP metod je to kernel spolu s hodnotou nebo rozdělením hyperparametrů, u PFN je to implicitní prior daný trénovacím rozdělením.
V tomto experimentu tedy nejde o špatně odhadnutou korelační délku, kterou by šlo z dat opravit.
Špatně je samotná třída funkcí, Matérnovy realizace totiž žádná volba hyperparametrů RBF přesně nepopíše.
\par

Jinou roli zde hraje i velikost kontextu.
V sekcích~\ref{sec:vnitrni-struktura} a~\ref{sec:pfn-random-hp} bylo vše správně specifikované a $n$ měřilo, jak rychle PFN konverguje ke správnému posterioru.
Zde je prior všech metod špatně a $n$ měří, jak rychle se metody ze špatného prioru vzpamatují.
S rostoucím $n$ totiž likelihood postupně převáží prior.
Adaptivní metody GP\_ML a GP\_Bayes si korelační délku mohou z dat přeodhadnout, jejich handicap tedy s daty mizí.
PFN svůj prior za běhu změnit nemůže a GP\_fixed se nepřizpůsobuje vůbec.
Nejzajímavější proto budou malé kontexty, kde prior dominuje a ukazuje se, čí prior škodí nejméně.
\par

\subsection{Vliv velikosti kontextu}
\label{subsec:exp7-q1}

V prvním testu zafixujeme $\ell_{\text{true}} = 0{,}3$ a měníme velikost kontextu $n \in \{5, 10, 20, 40, 64, 100, 128\}$.
Ptáme se, jak rychle se jednotlivé metody s rostoucím množstvím dat vzpamatují ze špatného prioru.
Výsledky shrnují tabulky~\ref{tab:exp7_mse_n} a~\ref{tab:exp7_dnll_n} a obrázek~\ref{fig:exp7_q1}.
Horní panel obrázku ukazuje MSE jako funkci $n$, dolní panel $\Delta$NLL.
Tečkovaná čára v dolním panelu značí \emph{oracle}, tedy nulu, ke které se chtějí všechny metody přiblížit.

\begin{figure}[p]
    \centering
    \begin{minipage}{\linewidth}
    \centering
    \begin{tabular}{cccccc}
    \hline
    $n$ & PFN & GP\_fixed & GP\_ML & GP\_Bayes & GP\_oracle \\
    \hline
    $5$   & $0{,}93699$ & $1{,}18172$ & $0{,}95305$ & $0{,}98592$ & $0{,}79025$ \\
    $10$  & $0{,}72713$ & $1{,}00189$ & $0{,}72144$ & $0{,}73947$ & $0{,}66794$ \\
    $20$  & $0{,}52432$ & $0{,}78485$ & $0{,}55354$ & $0{,}54532$ & $0{,}47545$ \\
    $40$  & $0{,}37256$ & $0{,}62423$ & $0{,}41614$ & $0{,}39967$ & $0{,}35816$ \\
    $64$  & $0{,}25835$ & $0{,}49483$ & $0{,}26753$ & $0{,}26891$ & $0{,}24434$ \\
    $100$ & $0{,}18719$ & $0{,}38220$ & $0{,}18626$ & $0{,}18506$ & $0{,}17135$ \\
    $128$ & $0{,}17385$ & $0{,}41661$ & $0{,}17169$ & $0{,}18200$ & $0{,}16918$ \\
    \hline
    \end{tabular}
    \captionof{table}{MSE vůči skutečným hodnotám jako funkce velikosti kontextu $n$ při misspecifikaci kernelu ($\ell_{\text{true}} = 0{,}3$, průměr přes $150$ instancí). PFN je pro $n \in \{10, 20, 40\}$ nejlepší ze všech misspecifikovaných metod, pro velká $n$ je s GP\_ML a GP\_Bayes prakticky shodné.}
    \label{tab:exp7_mse_n}
    \end{minipage}
    \vfill
    \begin{minipage}{\linewidth}
    \centering
    \begin{tabular}{ccccc}
    \hline
    $n$ & PFN & GP\_fixed & GP\_ML & GP\_Bayes \\
    \hline
    $5$   & $+0{,}165$ & $+0{,}732$ & $+0{,}968$ & $+0{,}326$ \\
    $10$  & $+0{,}143$ & $+0{,}956$ & $+0{,}403$ & $+0{,}288$ \\
    $20$  & $+0{,}086$ & $+1{,}068$ & $+0{,}239$ & $+0{,}190$ \\
    $40$  & $+0{,}077$ & $+1{,}231$ & $+0{,}191$ & $+0{,}169$ \\
    $64$  & $+0{,}077$ & $+1{,}157$ & $+0{,}050$ & $+0{,}086$ \\
    $100$ & $+0{,}045$ & $+0{,}972$ & $+0{,}056$ & $+0{,}049$ \\
    $128$ & $+0{,}031$ & $+1{,}174$ & $+0{,}016$ & $+0{,}044$ \\
    \hline
    \end{tabular}
    \captionof{table}{Kalibrace $\Delta$NLL vůči \emph{oracle} jako funkce velikosti kontextu $n$ ($\ell_{\text{true}} = 0{,}3$, průměr přes $150$ instancí, \emph{oracle} má z definice hodnotu $0$). PFN je nejblíže \emph{oracle} pro všechna $n \leq 40$, teprve od $n = 64$ jej mírně porazí GP\_ML.}
    \label{tab:exp7_dnll_n}
    \end{minipage}
    \vfill
    \includegraphics[width=0.576\linewidth]{figures/GP2_exp7_misspecification/fig_exp7_q1_mse_vs_n.png}
    \caption{Vliv velikosti kontextu při misspecifikaci kernelu (data z Matérn $2{,}5$, $\ell_{\text{true}} = 0{,}3$). Nahoře: MSE jako funkce $n$, černá tečkovaná čára je \emph{GP\_oracle}. Dole: $\Delta$NLL vůči \emph{oracle}, tečkovaná čára značí nulu (\emph{oracle}). Chybové úsečky $\pm 1$ směrodatná chyba. PFN (modrá) se v MSE drží těsně u \emph{oracle} a v $\Delta$NLL je pro malé a střední kontexty ze všech metod nejníže.}
    \label{fig:exp7_q1}
\end{figure}

\textbf{Diskuze:}

Z tabulek plynou dvě pozorování.
Za prvé, v MSE je PFN plně konkurenceschopné s adaptivními metodami GP\_ML a GP\_Bayes.
Pro střední velikosti kontextu je dokonce nejlepší ze všech misspecifikovaných metod a pro velká $n$ jsou všechny tři prakticky shodné, jak je vidět v tabulce~\ref{tab:exp7_mse_n}.
GP\_fixed, který se datům nijak nepřizpůsobuje, zaostává za všemi výrazně a s rostoucím $n$ se jeho odstup nezmenšuje.
Za druhé, v kalibraci je obraz ještě jednoznačnější.
PFN má nejmenší $\Delta$NLL ze všech misspecifikovaných metod pro všechna $n \leq 40$ a teprve pro velké kontexty jej GP\_ML porazí, jak ukazuje tabulka~\ref{tab:exp7_dnll_n}.
\par

Toto chování má přirozené vysvětlení.
Adaptivní metody musí korelační délku odhadnout z několika málo bodů a při malém kontextu ji přefitují, jejich prediktivní intervaly jsou pak příliš sebejisté.
PFN naproti tomu marginalizuje přes celé trénovací rozdělení hyperparametrů, a proto vrací opatrnější a lépe kalibrovanou nejistotu.
Přesně v režimu malých dat, kde na kalibraci záleží nejvíce, se tato vlastnost vyplácí.
\par

\subsection{Vliv síly misspecifikace}
\label{subsec:exp7-q2}

Ve druhém testu zafixujeme $n = 40$ a měníme pravou korelační délku $\ell_{\text{true}} \in [0{,}05,\, 2{,}0]$.
Tím měníme sílu misspecifikace.
Pro malá $\ell_{\text{true}}$ jsou realizace Matérnova kernelu hrubé a RBF je aproximuje špatně.
Pro velká $\ell_{\text{true}}$ jsou realizace hladké a rozdíl mezi oběma kernely se stírá.
Hodnoty $\ell_{\text{true}} > 1{,}0$ navíc leží mimo trénovací rozsah PFN, takže zároveň testujeme, jak model degraduje mimo svůj implicitní prior.
Výsledky shrnují tabulky~\ref{tab:exp7_mse_ls} a~\ref{tab:exp7_dnll_ls} a obrázek~\ref{fig:exp7_q2}.

\begin{table}[htbp]
\centering
\begin{tabular}{cccccc}
\hline
$\ell_{\text{true}}$ & PFN & GP\_fixed & GP\_ML & GP\_Bayes & GP\_oracle \\
\hline
$0{,}05$ & $1{,}06612$ & $2{,}00588$ & $0{,}99193$ & $0{,}98868$ & $0{,}98116$ \\
$0{,}10$ & $0{,}85773$ & $1{,}87695$ & $0{,}81220$ & $0{,}81922$ & $0{,}78140$ \\
$0{,}20$ & $0{,}55296$ & $1{,}13399$ & $0{,}57635$ & $0{,}56389$ & $0{,}50696$ \\
$0{,}30$ & $0{,}37256$ & $0{,}62423$ & $0{,}41614$ & $0{,}39967$ & $0{,}35816$ \\
$0{,}50$ & $0{,}24395$ & $0{,}27506$ & $0{,}25383$ & $0{,}23612$ & $0{,}23189$ \\
$0{,}70$ & $0{,}19586$ & $0{,}18896$ & $0{,}19154$ & $0{,}18329$ & $0{,}18401$ \\
$1{,}00$ & $0{,}16668$ & $0{,}15438$ & $0{,}15891$ & $0{,}15302$ & $0{,}15353$ \\
$1{,}50$ & $0{,}15210$ & $0{,}14165$ & $0{,}13702$ & $0{,}13590$ & $0{,}13366$ \\
$2{,}00$ & $0{,}14713$ & $0{,}13830$ & $0{,}13038$ & $0{,}12932$ & $0{,}12545$ \\
\hline
\end{tabular}
\caption{MSE vůči skutečným hodnotám jako funkce pravé korelační délky $\ell_{\text{true}}$ ($n = 40$, průměr přes $150$ instancí). Pro krátké korelační délky vede PFN spolu s adaptivními metodami, pro dlouhé se všechny metody sbíhají a GP\_fixed přestává být špatný, jeho $\ell = 0{,}7$ se totiž blíží pravdě.}
\label{tab:exp7_mse_ls}
\end{table}

\begin{figure}[p]
    \centering
    \begin{minipage}{\linewidth}
    \centering
    \begin{tabular}{ccccc}
    \hline
    $\ell_{\text{true}}$ & PFN & GP\_fixed & GP\_ML & GP\_Bayes \\
    \hline
    $0{,}05$ & $+0{,}083$ & $+5{,}544$ & $+0{,}007$ & $+0{,}007$ \\
    $0{,}10$ & $+0{,}098$ & $+5{,}078$ & $+0{,}041$ & $+0{,}058$ \\
    $0{,}20$ & $+0{,}099$ & $+2{,}758$ & $+0{,}148$ & $+0{,}145$ \\
    $0{,}30$ & $+0{,}077$ & $+1{,}231$ & $+0{,}191$ & $+0{,}169$ \\
    $0{,}50$ & $+0{,}057$ & $+0{,}241$ & $+0{,}098$ & $+0{,}058$ \\
    $0{,}70$ & $+0{,}055$ & $+0{,}044$ & $+0{,}052$ & $+0{,}022$ \\
    $1{,}00$ & $+0{,}061$ & $+0{,}005$ & $+0{,}034$ & $+0{,}006$ \\
    $1{,}50$ & $+0{,}085$ & $+0{,}025$ & $+0{,}015$ & $+0{,}007$ \\
    $2{,}00$ & $+0{,}097$ & $+0{,}043$ & $+0{,}020$ & $+0{,}013$ \\
    \hline
    \end{tabular}
    \captionof{table}{Kalibrace $\Delta$NLL vůči \emph{oracle} jako funkce $\ell_{\text{true}}$ ($n = 40$, průměr přes $150$ instancí). Pro střední korelační délky kolem pravé hodnoty $0{,}2$ až $0{,}3$ je nejblíže \emph{oracle} PFN. Pro velmi krátké i dlouhé délky vedou adaptivní metody, $\Delta$NLL PFN však zůstává malé v celém rozsahu, a to i mimo jeho trénovací rozsah $\ell > 1{,}0$.}
    \label{tab:exp7_dnll_ls}
    \end{minipage}
    \vfill
    \includegraphics[width=0.64\linewidth]{figures/GP2_exp7_misspecification/fig_exp7_q2_mse_vs_ls.png}
    \caption{Vliv síly misspecifikace ($n = 40$). Nahoře: MSE jako funkce $\ell_{\text{true}}$. Dole: $\Delta$NLL vůči \emph{oracle}. Modrý pás značí trénovací rozsah PFN $\ell \in [0{,}05,\, 1{,}0]$. GP\_fixed (oranžová) pro krátké korelační délky selhává v obou metrikách, ostatní metody se s rostoucím $\ell_{\text{true}}$ sbíhají k \emph{oracle}. PFN zůstává blízko \emph{oracle} v celém rozsahu včetně hodnot mimo svůj trénovací rozsah.}
    \label{fig:exp7_q2}
\end{figure}

\textbf{Diskuze:}

Za prvé, síla misspecifikace se projevuje přesně tam, kde bychom čekali.
Pro krátké korelační délky, kde jsou Matérnovy realizace hrubé, selhává GP\_fixed v kalibraci o řády, jak je vidět v tabulce~\ref{tab:exp7_dnll_ls}.
Pro dlouhé korelační délky jsou realizace hladké, RBF s $\ell = 0{,}7$ se stává rozumnou volbou a všechny metody se sbíhají.
Za druhé, PFN je nejlépe kalibrované v okolí pravé hodnoty $\ell_{\text{true}} \approx 0{,}3$, tedy uprostřed svého trénovacího rozsahu.
Pro extrémně krátké délky ji přebijí ostatní adaptivní metody, kterým hrubá data poskytují silný signál pro odhad $\ell$.
Za třetí, mimo trénovací rozsah ($\ell_{\text{true}} > 1{,}0$) PFN degraduje jen zlehka.
Jeho $\Delta$NLL zůstává malé i pro dvojnásobek horní hranice trénovacího rozdělení, žádný náhlý kolaps nenastává.
To je stejná plynulá degradace, jakou jsme u identifikace hyperparametrů pozorovali v sekci~\ref{sec:pfn-random-hp}.
\par

\subsection{Chování mimo trénovací rozsah korelační délky}
\label{subsec:exp7-ood}

Předchozí část naznačovala, že PFN může relativně dobře zvládnout i odhady pro korelační délky mimo svůj trénovací rozsah.
Avšak pozorný čtenář by mohl hned oponovat řadou argumentů.
Uvedený test totiž pro takový závěr není korektní.
Jedním z důvodů, který leží na povrchu, je tento.
Ačkoli některé hodnoty leží mimo rozsah, jsou to hodnoty na hladkém konci škály pro parametr $\ell$, kde je úloha sama o sobě snadná.
Protože odhadovat hladké funkce není těžká úloha a případné selhání by se skoro neprojevilo.
Čistý test vyžaduje dvě věci.
Kernel musí být správně specifikovaný, aby jediným zdrojem chyby byla extrapolace.
Dále chtěli bychom, aby korelační délka opustila trénovací rozsah na obou stranách.
Přesně takový test nyní provedeme.
\par

Data generujeme z GP s RBF kernelem, tedy ze stejné třídy funkcí, na které byl model trénován.
Pravou korelační délku měníme na logaritmické mřížce $\ell \in [0{,}01,\, 5{,}0]$, která přesahuje trénovací rozsah $[0{,}05,\, 1{,}0]$ na obou stranách.
Použijeme šestivrstvý náhodný model, $n = 40$ kontextových bodů, šum $\sigma^2 = 0{,}01$ a $200$ instancí pro každou hodnotu $\ell$.
Chybu měříme jako \emph{relativní MSE}:
$$\text{rel\_MSE} = \frac{\text{MSE}(\hat\mu_{\text{PFN}},\, \mu_{\text{oracle}})}{\text{MSE}(\mu_{\text{oracle}},\, \mathbf{y}_{\text{test}})},$$
tedy odchylku PFN od \emph{oracle} vydělenou chybou, kterou má sám \emph{oracle} vůči skutečným zašuměným hodnotám.
Jmenovatel je tedy chyba, které se nedá zbavit, neboť je způsobena samotným šumem v datech, a žádná metoda se pod ni nedostane.
Hodnota $\text{rel\_MSE} = 1$ by tedy znamenala, že se PFN liší od \emph{oracle} o tolik, kolik činí tato nevyhnutelná chyba.
Hodnoty hluboko pod $1$ by znamenaly, že PFN \emph{oracle} prakticky kopíruje.
Výsledky shrnuje tabulka~\ref{tab:exp7_ood} a obrázek~\ref{fig:exp7_ood}.

\begin{figure}[!htb]
    \centering
    \begin{minipage}{\linewidth}
    \centering
\begin{tabular}{ccccl}
\hline
$\ell$ & $\text{MSE}(\hat\mu_{\text{PFN}},\, \mu_{\text{oracle}})$ & $\text{MSE}(\mu_{\text{oracle}},\, \mathbf{y}_{\text{test}})$ & $\text{rel\_MSE}$ & Stav \\
\hline
$0{,}01$ & $0{,}35358$ & $0{,}46192$ & $0{,}765$ & mimo rozsah \\
$0{,}03$ & $0{,}07729$ & $0{,}09744$ & $0{,}793$ & mimo rozsah \\
$0{,}05$ & $0{,}01634$ & $0{,}03569$ & $0{,}458$ & hranice \\
$0{,}10$ & $0{,}00389$ & $0{,}01583$ & $0{,}246$ & \\
$0{,}20$ & $0{,}00198$ & $0{,}01233$ & $0{,}161$ & \\
$0{,}30$ & $0{,}00117$ & $0{,}01154$ & $0{,}101$ & minimum \\
$0{,}50$ & $0{,}00244$ & $0{,}01091$ & $0{,}224$ & \\
$0{,}70$ & $0{,}00318$ & $0{,}01069$ & $0{,}297$ & \\
$1{,}00$ & $0{,}00356$ & $0{,}01053$ & $0{,}338$ & hranice \\
$1{,}50$ & $0{,}00294$ & $0{,}01043$ & $0{,}282$ & mimo rozsah \\
$2{,}00$ & $0{,}00234$ & $0{,}01035$ & $0{,}226$ & mimo rozsah \\
$3{,}00$ & $0{,}00197$ & $0{,}01028$ & $0{,}191$ & mimo rozsah \\
$5{,}00$ & $0{,}00192$ & $0{,}01024$ & $0{,}188$ & mimo rozsah \\
\hline
\end{tabular}
\captionof{table}{Chování PFN mimo trénovací rozsah korelační délky (RBF data, $n = 40$, $\sigma^2 = 0{,}01$, průměr přes $200$ instancí). Relativní MSE je odchylka PFN od \emph{oracle} dělená chybou, kterou má sám \emph{oracle} kvůli šumu v datech. Minimum leží uprostřed trénovacího rozsahu a ani mimo rozsah na žádné straně relativní MSE nepřekročí referenční hodnotu $1$.}
\label{tab:exp7_ood}
    \end{minipage}
    \vspace{1.5em}

    \includegraphics[width=0.92\linewidth]{figures/GP2_exp6_hp_identification/fig_exp6_q4_ood.png}
    \caption{Relativní MSE jako funkce pravé korelační délky $\ell$ (logaritmická osa). Zelený pás značí trénovací rozsah $\ell \in [0{,}05,\, 1{,}0]$, zelená tečkovaná čára referenční hodnotu $\text{rel\_MSE} = 1$. Křivka má minimum uprostřed trénovacího rozsahu a mimo něj na obou stranách roste jen pozvolna, referenční hodnotu nikde nepřekročí.}
    \label{fig:exp7_ood}
\end{figure}

\textbf{Diskuze:}

Z tabulky~\ref{tab:exp7_ood} můžeme učinit jasné závěry.
Uvnitř trénovacího rozsahu je relativní MSE hluboko pod referenční hodnotou $\text{rel\_MSE} = 1$, tedy pod úrovní chyby způsobené samotným šumem v datech.
Nejmenší je uprostřed rozsahu.
Přesnou příčinu polohy minima neurčujeme, prior nad $\ell$ je rovnoměrný, takže minimum nelze vysvětlit tím, že by model střední délky viděl při tréninku častěji, jde spíše o~souhru obtížnosti úlohy a normalizace \emph{oraclem}.
Mimo rozsah chyba na obou stranách roste, ale referenční hodnotu nikde nepřekročí.
Odchylka PFN od \emph{oracle} tedy zůstává menší než chyba způsobená samotným šumem v datech.
A to pro případy, které jsou hodně vzdálené od hranice trénovací množiny, ale to opět může být způsobeno tím ,že pro tato $\ell$ je tato úloha výrazně jednodušší.
Z experimentu vyplývá, že k degradaci sice dochází, ale žádný náhlý kolaps nenastává.
\par

Z toho usuzujeme, že implicitní prior PFN nefunguje jako ostrá hranice, za kterou model přestává fungovat.
Chová se spíše jako měkké těžiště, od kterého přesnost pozvolna klesá.
Pro hladké funkce za horní hranicí rozsahu to dává dobrý smysl, model se během tréninku naučil reprezentovat hladkost a delší korelační délka pro něj není kvalitativně nová situace.
Pro velmi krátké délky pod dolní hranicí je úloha těžká i pro samotný \emph{oracle}, kontext $40$ bodů rychlé oscilace prostě nezachytí.
Tím je zpětně podložený i závěr předchozí části, plynulá degradace mimo trénovací rozsah není artefakt misspecifikovaného kernelu, platí i v čistém testu.
\par

\subsection{Shrnutí}
\label{subsec:exp7-shrnuti}

Tato sekce zkoumala otázku, co se stane, když správný prior nezná žádná z metod a jak se chová PFN v porovnání s nimi.
\par

Zjistili jsme, že ve smyslu MSE je PFN srovnatelné s adaptivními metodami GP\_ML a GP\_Bayes.
Pro střední velikosti kontextu je dokonce nejlepší, a to přesto, že na rozdíl od nich nemůže svůj prior za běhu nijak přizpůsobit datům.
V kalibraci je PFN nejlepší ze všech misspecifikovaných metod pro malý a střední kontext i pro korelační délky v okolí pravé hodnoty.
GP\_fixed, který nemá žádný mechanismus adaptace, selhává ve všech režimech nejvíce.
Čistý test s korektně specifikovaným kernelem navíc ukázal, že mimo trénovací rozsah korelační délky PFN degraduje plynule a jeho odchylka od \emph{oracle} zůstává menší než chyba způsobená samotným šumem v datech, jak je vidět v tabulce~\ref{tab:exp7_ood}.
\par

Vysvětlení je stejné jako u kalibrace rozptylu v sekci~\ref{sec:pfn-random-hp}.
PFN neprovádí bodový odhad hyperparametrů, ale marginalizuje přes celé trénovací rozdělení.
Tam, kde adaptivní metody z malého kontextu korelační délku přefitují a vracejí příliš sebejisté intervaly, drží PFN opatrnější nejistotu, a právě proto je jeho $\Delta$NLL nejmenší.
Za tuto robustnost platí jen malou cenu v MSE pro velké kontexty, kde už mají adaptivní metody dost dat na spolehlivý odhad.
\par

Dohromady tedy platí, že amortizovaná Bayesovská inference PFN není křehká vůči volbě prioru, alespoň v těchto syntetických experimentech.
I když se skutečný generativní proces liší od trénovacího prioru samotným kernelem, PFN zůstává přesné a ze všech srovnávaných metod nejlépe kalibrované právě v režimu malých dat.
